% fig_lec08_polar_radial.tex — Radial parallel transport in flat R² (polar coords)
% Data: generated by scripts/sim_lec07_lec08.jl
%
% A Cartesian-constant vector v = (1, 0) at every point along the radial ray
% φ = π/6.  In polar components,
%   v^r(r)  ≡  cos(φ_0) = √3/2  (constant)
%   v^φ(r)  =  −sin(φ_0)/r = −1/(2r)
% The components rotate along the ray purely because the polar basis
% (∂/∂r, ∂/∂φ) itself rotates; the vector itself is parallel.

\begin{center}
\begin{tikzpicture}
  % ── LEFT: Cartesian / polar plane with transported vector ───
  \begin{axis}[
      name=plane,
      width=6.4cm, height=6.4cm,
      axis equal image,
      axis lines=middle,
      axis line style={birrenorange, line width=0.6pt,
                       -{Stealth[length=5pt,width=3pt]}},
      xmin=-1.0, xmax=3.4, ymin=-0.8, ymax=2.4,
      xtick=\empty, ytick=\empty, clip=false,
  ]
    % Concentric grid circles (six datasets in one file separated by blank lines)
    \addplot[spacecadet!25, line width=0.4pt, no markers, forget plot]
      table {data/lec08_polar_circles.dat};

    % Radial ray φ = π/6
    \addplot[spacecadet!50, line width=0.6pt, no markers, forget plot,
            dashed]
      table {data/lec08_polar_ray.dat};

    % Transported vectors at sample radii (Cartesian-constant, drawn from r·(cos,sin)φ_0)
    \addplot[
      quiver={u=\thisrowno{2}, v=\thisrowno{3}},
      -{Stealth[length=5pt, width=3.5pt]},
      banana!70!black, line width=1.0pt, no markers, forget plot,
    ] table[x index=0, y index=1] {data/lec08_polar_vectors.dat};

    % Sample points
    \addplot[only marks, mark=*, mark size=1.5pt, banana!90!black, forget plot]
      table {data/lec08_polar_points.dat};

    % Labels
    \node[font=\sf\small, text=birrenorange] at (axis cs:3.25,2.20) {$\Rn{2}$};
    \node[font=\scriptsize, text=spacecadet!70, anchor=west]
      at (axis cs:2.92,1.46) {$\varphi=\tfrac{\pi}{6}$};
    \node[font=\scriptsize, text=banana!50!black, anchor=south]
      at (axis cs:1.0,0.82) {$v$};
  \end{axis}

  % ── RIGHT: components v^r, v^φ as functions of r ────────────
  \begin{axis}[
      name=comp,
      at={($(plane.east)+(1.6cm,0)$)}, anchor=west,
      width=5.4cm, height=4.6cm,
      axis lines=left,
      axis line style={birrenorange, line width=0.6pt,
                       -{Stealth[length=5pt,width=3pt]}},
      xmin=0.3, xmax=3.0, ymin=-1.7, ymax=1.1,
      xlabel={$r$}, x label style={at={(axis description cs:1.0,0)}, anchor=west,
                                   font=\scriptsize, text=spacecadet},
      tick label style={font=\tiny, text=spacecadet},
      xtick={1, 2, 3}, ytick={-1.5, -1, -0.5, 0, 0.5, 1},
      yticklabel style={/pgf/number format/fixed},
      clip=false,
  ]
    % v^r = cos(π/6)
    \addplot[banana!70!black, line width=1.0pt, no markers, forget plot]
      table {data/lec08_polar_vr.dat};
    % v^φ = -sin(π/6)/r
    \addplot[birrenred!85!black, line width=1.0pt, no markers, forget plot]
      table {data/lec08_polar_vphi.dat};

    \node[font=\scriptsize, text=banana!55!black, anchor=west]
      at (axis cs:2.40,0.95) {$v^r=\cos\varphi_0$};
    \node[font=\scriptsize, text=birrenred!85!black, anchor=west]
      at (axis cs:1.30,-0.55) {$v^\varphi=A/r$};
  \end{axis}
\end{tikzpicture}
\end{center}
